%% Official CVPR 2026 author-kit preamble (CVPR2026-v1), plus the extra packages
%% this paper needs. Everything above the "PAPER-SPECIFIC" rule is the stock kit
%% content, left untouched so it stays easy to diff against a future kit release.
%%
%% This file contains a number of tweaks that are typically applied to the main document.
%% They are not enabled by default, but can be enabled by uncommenting the relevant lines.

%%
%% Inline annotations; for predefined colors, refer to "dvipsnames" in the xcolor package:
%% https://tinyurl.com/overleaf-colors
%%
\newcommand{\red}[1]{{\color{red}#1}}
\newcommand{\todo}[1]{{\color{red}#1}}
\newcommand{\TODO}[1]{\textbf{\color{red}[TODO: #1]}}
%%
%% disable for camera ready / submission by uncommenting these lines
%%
% \renewcommand{\TODO}[1]{}
% \renewcommand{\todo}[1]{#1}

%%
%% work harder in optimizing text layout. Typically shrinks text by 1/6 of page, enable
%% it at the very end of the writing process, when you are just above the page limit
%%
% \usepackage{microtype}

%%
%% fine-tune paragraph spacing
%%
% \renewcommand{\paragraph}[1]{\vspace{.5em}\noindent\textbf{#1.}}

%%
%% globally adjusts space between figure and caption
%%
% \setlength{\abovecaptionskip}{.5em}


%%
%% Allows "the use of \paper to refer to the project name"
%% with automatic management of space at the end of the word
%%
% \usepackage{xspace}
% \newcommand{\paper}{ProjectName\xspace}

%%
%% Commonly used math definitions
%%
% \DeclareMathOperator*{\argmin}{arg\,min}
% \DeclareMathOperator*{\argmax}{arg\,max}

%%
%% Tigthen underline
%%
% \usepackage{soul}
% \setuldepth{foobar}

% =====================================================================
% PAPER-SPECIFIC ADDITIONS (not part of the stock author kit)
% =====================================================================
% cvpr.sty already provides: times, xspace, xcolor, graphicx, amsmath, amssymb,
% booktabs, natbib, caption/subcaption, enumitem, and it loads cleveref itself
% via \AtEndPreamble (i.e. after hyperref). Only the following are missing.
\usepackage{algorithm}
\usepackage{algpseudocode}

% Figures are pre-converted PDFs, so the svg package is NOT needed. Uncomment it
% only if you switch a figure back to \includesvg[inkscapelatex=false]{...}.
% \usepackage{svg}

% ---------------------------------------------------------------------
% Review line numbers: keep them out of the text.
%
% cvpr.sty loads lineno with [switch], which picks the side from "am I in the
% first or second column?". In a two-column document with full-width floats
% that test is unreliable: here it came out inverted, so BOTH columns' numbers
% were set into the 22pt gutter. With the stock \linenumbersep of .75cm (21pt)
% plus a 13pt-wide three-digit number, each one then ran into the neighbouring
% column's text (~500 collisions over 9 pages).
%
%   \switchlinenumbers*  -- flip the side test back, so numbers land in the
%                           OUTER page margins (left column -> left margin,
%                           right column -> right margin).
%   \linenumbersep 8pt   -- belt and braces: small enough that a number still
%                           fits inside the gutter on any page where lineno
%                           guesses "inner" anyway (page 1, whose column sense
%                           is thrown off by \twocolumn[\maketitle]).
%
% This is a deviation from the stock author kit; it changes only where the
% numbers sit, never the numbering itself. Guarded so the camera-ready build
% (no [review] option => no lineno) is unaffected.
% ---------------------------------------------------------------------
\makeatletter
\@ifpackageloaded{lineno}{%
  \switchlinenumbers*%
  \setlength{\linenumbersep}{8pt}%
}{}
\makeatother
